% GENERATED FILE. Edit canonical lessons, then run npm run generate:books.
% canonical-source-hash: fnv1a64:a5a41c0b3e45a7ab
% canonical-lessons: JA-C172-shokuji, JA-C172-hirugohan, JA-C172-suku, JA-C172-bangohan, JA-C172-yaku

\chapter{Mealtimes and Hunger}
\label{ch:ja-mealtimes-and-hunger}
\hlchaptermodality{172}

\begin{chapteropening}
\textbf{By the end of this chapter:} \emph{I can talk about my meals and when I eat them.}

Complete the last lesson of chapter 172: I can talk about my meals and when I eat them.
\end{chapteropening}

\section[\texorpdfstring{\ja{しょくじ} (shokuji)}{しょくじ (shokuji)}]{\ja{しょくじ} (shokuji) — meal; dining}

\label{lesson:JA-C172-shokuji}



\begin{quote}
Before the new one: say the Japanese for open for business, then the Japanese for not yet open.
\end{quote}

\subsection*{You'll want to know: \ja{しょくじ}}
\textbf{\ja{しょくじ}} — \emph{shokuji} — \textquotedblleft{}meal; dining\textquotedblright{}.

\begin{quote}
\textbf{\ja{まいにち} \ja{かぞくと} \ja{しょくじを} \ja{します。}} — \emph{mainichi kazoku to shokuji o shimasu} — \textquotedblleft{}I have meals with my family every day.\textquotedblright{}
\end{quote}

\begin{cousinweb}
\textbf{\ja{しょくじ}} (\emph{shokuji}) is Sino-Japanese (kango): two Chinese characters read with their on'yomi, readings borrowed from Chinese. They are \emph{shoku}, \textquotedblleft{}eat, food\textquotedblright{} (Middle Chinese \emph{zyik}) and \emph{ji}, \textquotedblleft{}matter, thing\textquotedblright{} (\emph{dzriH}) (KANJIDIC2; Middle Chinese in Baxter's transcription).
\end{cousinweb}

\begin{grammarlens}[title={What you've built}]
One more word for everyday life at A2.
\end{grammarlens}

\subsection*{Your turn}
\begin{itemize}
\raggedright
  \item \emph{Say it:} \emph{shokuji}
  \item \emph{Say it:} \emph{shokuji}, once more
  \item \emph{Say it:} \emph{mainichi kazoku to shokuji o shimasu} — I have meals with my family every day.
  \item \emph{Recall:} say the Japanese for open for business, then the Japanese for not yet open, then say \emph{shokuji} again
\end{itemize}

\begin{tcolorbox}[breakable,colback=teal!4,colframe=teal!35!black,title={Before you move on}]
What does \ja{しょくじ} mean? (\textquotedblleft{}Meal; dining\textquotedblright{}.) Say it once more.
\end{tcolorbox}
\section[\texorpdfstring{\ja{ひるごはん} (hirugohan)}{ひるごはん (hirugohan)}]{\ja{ひるごはん} (hirugohan) — lunch}

\label{lesson:JA-C172-hirugohan}



\begin{quote}
Before the new one: say the Japanese for not yet open, then the Japanese for meal.
\end{quote}

\subsection*{You'll want to know: \ja{ひるごはん}}
\textbf{\ja{ひるごはん}} — \emph{hirugohan} — \textquotedblleft{}lunch\textquotedblright{}.

\begin{quote}
\textbf{\ja{ひるごはんに} \ja{うどんを} \ja{たべます。}} — \emph{hirugohan ni udon o tabemasu} — \textquotedblleft{}I eat udon for lunch.\textquotedblright{}
\end{quote}

\begin{cousinweb}
\textbf{\ja{ひるごはん}} (\emph{hirugohan}) is built from native \emph{hiru}, \textquotedblleft{}midday\textquotedblright{}, and Sino-Japanese \emph{gohan}, \textquotedblleft{}cooked rice, a meal\textquotedblright{}, whose two parts are readings borrowed from Chinese (JMdict, KANJIDIC2).
\end{cousinweb}

\begin{grammarlens}[title={What you've built}]
One more word for everyday life at A2.
\end{grammarlens}

\subsection*{Your turn}
\begin{itemize}
\raggedright
  \item \emph{Say it:} \emph{hirugohan}
  \item \emph{Say it:} \emph{hirugohan}, once more
  \item \emph{Say it:} \emph{hirugohan ni udon o tabemasu} — I eat udon for lunch.
  \item \emph{Recall:} say the Japanese for not yet open, then the Japanese for meal, then say \emph{hirugohan} again
\end{itemize}

\begin{tcolorbox}[breakable,colback=teal!4,colframe=teal!35!black,title={Before you move on}]
What does \ja{ひるごはん} mean? (\textquotedblleft{}Lunch\textquotedblright{}.) Say it once more.
\end{tcolorbox}
\section[\texorpdfstring{\ja{すく} (suku)}{すく (suku)}]{\ja{すく} (suku) — to get hungry; to become empty, to thin out}

\label{lesson:JA-C172-suku}



\begin{quote}
Before the new one: say the Japanese for meal, then the Japanese for lunch.
\end{quote}

\subsection*{You'll want to know: \ja{すく}}
\textbf{\ja{すく}} — \emph{suku} — \textquotedblleft{}to get hungry; to become empty, to thin out\textquotedblright{}.

\begin{quote}
\textbf{\ja{ゆうがた} \ja{おなかが} \ja{すきます。}} — \emph{yūgata onaka ga sukimasu} — \textquotedblleft{}I get hungry in the evening.\textquotedblright{}
\end{quote}

Polite \textbf{\ja{すきます}}, \emph{sukimasu}; a godan verb. The hungry one is the stomach: \textbf{\ja{おなかが} \ja{すきます}}.

\begin{cousinweb}
\textbf{\ja{すく}} (\emph{suku}) is a native Japanese word (wago), inherited rather than borrowed from Chinese; its kanji was attached later as a kun'yomi (JMdict, KANJIDIC2).
\end{cousinweb}

\begin{grammarlens}[title={What you've built}]
One more word for everyday life at A2.
\end{grammarlens}

\subsection*{Your turn}
\begin{itemize}
\raggedright
  \item \emph{Say it:} \emph{suku}
  \item \emph{Say it:} \emph{suku}, once more
  \item \emph{Say it:} \emph{yūgata onaka ga sukimasu} — I get hungry in the evening.
  \item \emph{Recall:} say the Japanese for meal, then the Japanese for lunch, then say \emph{suku} again
\end{itemize}

\begin{tcolorbox}[breakable,colback=teal!4,colframe=teal!35!black,title={Before you move on}]
What does \ja{すく} mean? (\textquotedblleft{}To get hungry; to become empty, to thin out\textquotedblright{}.) Say it once more.
\end{tcolorbox}
\section[\texorpdfstring{\ja{ばんごはん} (bangohan)}{ばんごはん (bangohan)}]{\ja{ばんごはん} (bangohan) — dinner; evening meal}

\label{lesson:JA-C172-bangohan}



\begin{quote}
Before the new one: say the Japanese for lunch, then the Japanese for to get hungry.
\end{quote}

\subsection*{You'll want to know: \ja{ばんごはん}}
\textbf{\ja{ばんごはん}} — \emph{bangohan} — \textquotedblleft{}dinner; evening meal\textquotedblright{}.

\begin{quote}
\textbf{\ja{ばんごはんは} \ja{いつも} \ja{うちで} \ja{たべます。}} — \emph{bangohan wa itsumo uchi de tabemasu} — \textquotedblleft{}I always eat dinner at home.\textquotedblright{}
\end{quote}

\begin{cousinweb}
\textbf{\ja{ばんごはん}} (\emph{bangohan}) is Sino-Japanese (kango): three Chinese characters read with their on'yomi, readings borrowed from Chinese. They are \emph{ban}, \textquotedblleft{}evening\textquotedblright{} (Middle Chinese \emph{mjonX}), \emph{go}, \textquotedblleft{}honourable\textquotedblright{} (\emph{ngjoH}) and \emph{han}, \textquotedblleft{}cooked rice\textquotedblright{} (\emph{bjonH}) (KANJIDIC2; Middle Chinese in Baxter's transcription).
\end{cousinweb}

\begin{grammarlens}[title={What you've built}]
One more word for everyday life at A2.
\end{grammarlens}

\subsection*{Your turn}
\begin{itemize}
\raggedright
  \item \emph{Say it:} \emph{bangohan}
  \item \emph{Say it:} \emph{bangohan}, once more
  \item \emph{Say it:} \emph{bangohan wa itsumo uchi de tabemasu} — I always eat dinner at home.
  \item \emph{Recall:} say the Japanese for lunch, then the Japanese for to get hungry, then say \emph{bangohan} again
\end{itemize}

\begin{tcolorbox}[breakable,colback=teal!4,colframe=teal!35!black,title={Before you move on}]
What does \ja{ばんごはん} mean? (\textquotedblleft{}Dinner; evening meal\textquotedblright{}.) Say it once more.
\end{tcolorbox}
\section[\texorpdfstring{\ja{やく} (yaku)}{やく (yaku)}]{\ja{やく} (yaku) — to grill; to bake, to roast}

\label{lesson:JA-C172-yaku}



\begin{quote}
Before the new one: say the Japanese for to get hungry, then the Japanese for dinner.
\end{quote}

\subsection*{You'll want to know: \ja{やく}}
\textbf{\ja{やく}} — \emph{yaku} — \textquotedblleft{}to grill; to bake, to roast\textquotedblright{}.

\begin{quote}
\textbf{\ja{にわで} \ja{にくを} \ja{やきます。}} — \emph{niwa de niku o yakimasu} — \textquotedblleft{}We grill meat in the garden.\textquotedblright{}
\end{quote}

Polite \textbf{\ja{やきます}}, \emph{yakimasu}; a godan verb.

\begin{cousinweb}
\textbf{\ja{やく}} (\emph{yaku}) is a native Japanese word (wago), inherited rather than borrowed from Chinese; its kanji was attached later as a kun'yomi (JMdict, KANJIDIC2).
\end{cousinweb}

\begin{grammarlens}[title={What you've built}]
One more word for everyday life at A2.
\end{grammarlens}

\subsection*{Your turn}
\begin{itemize}
\raggedright
  \item \emph{Say it:} \emph{yaku}
  \item \emph{Say it:} \emph{yaku}, once more
  \item \emph{Say it:} \emph{niwa de niku o yakimasu} — We grill meat in the garden.
  \item \emph{Recall:} say the Japanese for to get hungry, then the Japanese for dinner, then say \emph{yaku} again
\end{itemize}

\begin{tcolorbox}[breakable,colback=teal!4,colframe=teal!35!black,title={Before you move on}]
What does \ja{やく} mean? (\textquotedblleft{}To grill; to bake, to roast\textquotedblright{}.) Say it once more.
\end{tcolorbox}
